\section{Physical records and an exact marked identity}
Write an outgoing section point as
\[
 q=R n_\alpha,\qquad v=\sqrt{1-p^2}\,n_\alpha+p\,t_\alpha,
 \qquad (\alpha,p)\in\T\times(-1,1).
\]
The collision map is $T_R$, its next physical flight is $\tau_R$, and
\begin{equation}\label{eq:flux}
 \dd\nu=\frac{\dd\alpha\dd p}{4\pi},\qquad
 \bar\tau_R=\frac{\sqrt3/2-\pi R^2}{2R},\qquad
 \dd\mu_R=\bar\tau_R^{-1}\dd\nu\dd s\quad(0\le s<\tau_R).
\end{equation}
These formulas follow by integrating the boundary flux $\cos\phi\dd s_{\partial Q}\dd\phi$: total section flux is $4\pi R$ and phase volume is $2\pi(\sqrt3/2-\pi R^2)$. Reflection preserves this flux. Thus $T_R$ preserves $\nu$.

Each next disk has a lattice label $d\ne0$. With $D=d-q$ and $L=D\cdot v$, its candidate entry time is
\begin{equation}\label{eq:entry}
 \tau_d=L-\sqrt{R^2-|D|^2+L^2}.
\end{equation}
The physical branch is the positive entry root with the smallest time, not an arbitrarily selected root. Its endpoint normal is $n'=(q+\tau_dv-d)/R$, and its outgoing velocity is $v'=v-2(v\cdot n')n'$. These expressions define the label, point, both velocities, and time on every regular branch. On a compact finite itinerary separated from tangencies and competing roots they depend analytically on $R$ and the initial coordinates, by the implicit function theorem. This is a finite-itinerary assertion, not a common norm for all iterates.

Distinct disks have distance at least $g_R=1-2R\ge3/50$. A straight outgoing flight cannot return to its initial convex disk. In particular there is no finite-time collision accumulation. This family has uniform finite horizon: a rational lattice direction has row spacing at most $\sqrt3/2$, so every line in that direction meets an open disk when $R>\sqrt3/4$; an irrational line is dense on the torus. Compactness then excludes arbitrarily long free segments at the smallest radius. Increasing the radius preserves this upper bound. Tangencies and their finite preimages have section measure zero.

Put $S_j\tau(y)=\sum_{i=0}^{j-1}\tau(T_R^iy)$. For $y$ interpreted as the first future outgoing collision and $u$ as its arrival time, the first $k$ future collisions have times $u+S_j\tau(y)$, $0\le j<k$. Their points, lattice labels, and incoming and outgoing velocities are obtained by iterating \eqref{eq:entry}. The final incomplete flight has duration $t-u-S_{k-1}\tau(y)$.

\begin{proposition}[Marked stationary record]\label{prop:marked}
For every bounded measurable functional $\Psi$ of this record and its terminal flight,
\begin{align}\label{eq:marked}
 \mathbb E_{\mu_R}[\Psi;N_t=k]
 =\frac1{\bar\tau_R}\int\dd\nu(y)\int_0^{\tau(T_R^{-1}y)}
 &\one_{\{u+S_{k-1}\tau(y)\le t<u+S_k\tau(y)\}}\nonumber\\[-2mm]
 &\hspace{-8mm}\cdot\Psi(\mathcal R_{k,t}(y,u))\dd u,\qquad k\ge1.
\end{align}
The initial phase point is $\Phi_R^{-u}y$, so the entire initial residual flight is retained as well. For $t=11/100$ there are at most two future collisions, and \eqref{eq:marked} with $k=1,2$, together with the zero-collision case, gives the full marked law without an independence assumption.
\end{proposition}
\begin{proof}
Under \eqref{eq:flux}, set $y=T_Rx$ and $u=\tau(x)-s$. Invariance of $\nu$ gives the displayed domain for $u$ with Jacobian one. The $k$th and $(k+1)$st future collision times are exactly the two endpoints in the indicator. This proves the identity for arbitrary bounded measurable $\Psi$, first for indicators and then by integration. The event of equality at an endpoint has zero suspension measure. Three collisions require at least two complete intercollision flights after the first event; their total duration is at least $6/50>11/100$.
\end{proof}

